\documentclass[a4paper,12pt]{article}
\usepackage[utf8]{inputenc}
\usepackage[english,russian]{babel}
\usepackage[T2A]{fontenc}

\usepackage[
  a4paper, mag=1000, includefoot,
  left=1.1cm, right=1.1cm, top=1.2cm, bottom=1.2cm, headsep=0.8cm, footskip=0.8cm
]{geometry}

\usepackage{amsmath}
\usepackage{amssymb}
\usepackage{times}
\usepackage{mathptmx}

\newcommand{\br}[1]{\left(#1\right)}

\IfFileExists{pscyr.sty}
{
  \usepackage{pscyr}
  \def\rmdefault{ftm}
  \def\sfdefault{ftx}
  \def\ttdefault{fer}
  \DeclareMathAlphabet{\mathbf}{OT1}{ftm}{bx}{it} % bx/it or bx/m
}

\mathsurround=0.1em
\clubpenalty=1000%
\widowpenalty=1000%
\brokenpenalty=2000%
\frenchspacing%
\tolerance=2500%
\hbadness=1500%
\vbadness=1500%
\doublehyphendemerits=50000%
\finalhyphendemerits=25000%
\adjdemerits=50000%


\begin{document}

\author{Коренев Иван}
\title{Метод Гаусса решения линейной системы с выбором главного элемента по всей матрице.}
\date{\today}
\maketitle

% \begin{center}
% \bf{Постановка задачи}
% \end{center}


\section{Постановка задачи}
Решается система $Ax = y$ с вещественной квадратной матрицей коээфициентов размера $n$, где $x, y \in \mathbb{R}^n$ - столбцы неизвестных и свободных членов, соответственно, с помощью блочного метода Гаусса с выбором главного элемента по всей матрице.

\section{Описание хранения данных}
Матрица хранится в одномерном массиве $A$ поблочно, сами блоки хранятся построчно. Для массива $A$ выделяется ровно $n^2$ элементов. Вектора правой части и неизвестных хранятся в одномерных массивах длины $n$. Координаты элемента матрицы $A$ в самой матрице называются глобальными, в блоке - локальными.

Все индексы и номер шага начинаются с $0$.
\subsection{Разбиение на блоки}
Пусть $n = m k + l, k = \lfloor \frac{n}{m} \rfloor, 0 \leq l < m$. При $l = 0$ блоки с номером $k$ отсутствуют. При $l > 0$ разбиение имеет вид:
\[
    A =
    \begin{pmatrix}
        A_{0, 0}^{m \times m} & \cdots & A_{0, k - 1}^{m \times m} & A_{0, k}^{m \times l}     \\
        \vdots & \ddots & \vdots & \vdots                    \\
        A_{k - 1, 0}^{m \times m} & \cdots & A_{k - 1, k - 1}^{m \times m} & A_{k - 1, k}^{m \times l} \\
        A_{k, 0}^{l \times m} & \cdots & A_{k, k - 1}^{l \times m}. & A_{k, k}^{l \times l}
    \end{pmatrix}.
\]

При $k = 0$ имеется только блок $A_{0, 0}^{l \times l}$.

Размер блока в одном направлении равен $\operatorname{size}(i) = \min(m, n - i m)$. Связь локальных и глобальных координат имеет вид:
\[
    \left(A_{i, j}^{\operatorname{size}(i) \times \operatorname{size}(j)} \right)_{p, q} = a_{i m + p, j m + q},
\]
\[
    0 \leq i, j < \left\lceil \frac{n}{m} \right\rceil, 0 \leq p < \operatorname{size}(i), 0 \leq q < \operatorname{size}(j).
\]

Элемент $(p, q)$ блока с номерами $(i, j)$ находится в массиве $A$ по индексу
\[
    m (i n + j \operatorname{size}(i)) + p \operatorname{size}(j) + q,
\]
с теми же множеством индексов.

Номер шага указывается в скобках справа сверху.

\section{Описание функций}
В функциях $i, j$ - номера блока, $W$ - блок с этими координатам, размеры блока равны $wr = \operatorname{size}(i)$, $wc = \operatorname{size}(j)$. Элемент с локальными координатами $(p, q)$ хранится в $W$ по индексу $p$, $wc + q$.

\begin{enumerate}
    \item size(i) возвращает фактический размер блока в одном направлении.
    \item set\_block(A, i, j, W, wr, wc) копирует $wr \times wc$ элементов из $W$ в $A$, начиная с индекса $i m n + j m \times wr$. Остальные блоки не изменяются.
    \item get\_block(A, i, j, W, wr, wc) копирует $wr \times wc$ элементов из $A$, начиная с того же индекса, в $W$. Матрица $A$ не изменяется.
\end{enumerate}

\section{Формулы метода}
\subsection{Прямой ход метода Гаусса}

На шаге $\alpha = 0, \dots, k$ выбирается невырожденных блок, у которого норма обратной матрицы минимальна. Таким образом, координаты главного блока $(p, q)$ находятся
\[
    (p,q) \in \underset
    {\substack{\alpha \le i,j < k \\ \text{неворожденна}}}
    {\arg\min}
    \left\| \left(\br{A_{i,j}^{m \times m}}^{(\alpha)}\right)^{-1} \right\|
\],
где норма блока берется как корень из суммы квадратов элементов, $\left|\left| B\right|\right| = \sqrt{\sum b_{ij}^2}$, и все матрица, из которых выбираем являются невырожденными. При существовании нескольких вариантов выбора блока берем с наименьшими координатам.

После выбора главного блока производим замену блочных строк $\alpha$ и $p$ и блочных столбцов $\alpha$ и $q$, запоминаем перестановку столбцов, также переставляем соответствующие блоки в столбце свободных членов $y$. Далее индексы идут после переименования

На $\alpha$ шаге будем иметь следующих вид матрицы
\[
    \begin{pmatrix}
        E^{m \times m} & \br{A_{0, 1}^{m \times m}}^{(\alpha)} & \br{A_{0, 2}^{m \times m}}^{(\alpha)} & \dots  & \br{A_{0, \alpha}^{m \times m}}^{(\alpha)}          & \dots  & \br{A_{0, k - 1}^{m \times m}}^{(\alpha)}          & \br{A_{0, k}^{m \times l}}^{(\alpha)}          \\
        0              & E^{m \times m}                        & \br{A_{1, 2}^{m \times m}}^{(\alpha)} & \dots  & \br{A_{1, \alpha}^{m \times m}}^{(\alpha)}          & \dots  & \br{A_{1, k - 1}^{m \times m}}^{(\alpha)}          & \br{A_{1, k}^{m \times l}}^{(\alpha)}          \\
        \vdots         & \vdots                                & \vdots                                & \ddots & \vdots                                              & \vdots & \vdots                                             & \vdots                                         \\
        0              & \dots                                 & 0                                     & \dots  & \br{A_{\alpha, \alpha}^{m \times m}}^{(\alpha)}     & \dots  & \br{A_{\alpha, k - 1}^{m \times m}}^{(\alpha)}     & \br{A_{\alpha, k}^{m \times l}}^{(\alpha)}     \\
        0              & \dots                                 & 0                                     & \dots  & \br{A_{\alpha + 1, \alpha}^{m \times m}}^{(\alpha)} & \dots  & \br{A_{\alpha + 1, k - 1}^{m \times m}}^{(\alpha)} & \br{A_{\alpha + 1, k}^{m \times l}}^{(\alpha)} \\
        \vdots         & \vdots                                & \vdots                                & \ddots & \vdots                                              & \vdots & \vdots                                             & \vdots                                         \\
        0              & \dots                                 & 0                                     & \dots  & \br{A_{k - 1, \alpha}^{m \times m}}^{(\alpha)}      & \dots  & \br{A_{k - 1, k - 1}^{m \times m}}^{(\alpha)}      & \br{A_{k - 1, k}^{m \times l}}^{(\alpha)}      \\
        0              & \dots                                 & 0                                     & \dots  & \br{A_{k, \alpha}^{l \times m}}^{(\alpha)}          & \dots  & \br{A_{k, k - 1}^{l \times m}}^{(\alpha)}          & \br{A_{k, k}^{l \times l}}^{(\alpha)}
    \end{pmatrix}.
\]

Пусть $L = \br{\br{A_{\alpha, \alpha}^{m \times m}}^{(\alpha)}}^{-1}$, находим $L$ обычным методом Гаусса(операциями над строчками). Аналогично будем находить $R$, только операциями над стробцами. Приводим к единичному виду выбранный блок и изменяем соответствующую строку, делаем соответствующие преобразования со стробцов свободных членов.
\begin{align*}
    \br{A_{\alpha, \alpha}^{m \times m}}^{(\alpha + 1)}     & = E^{m \times m} \\
    \br{A_{\alpha, j}^{m \times m}}^{(\alpha + 1)}          & = L \br{A_{\alpha, j}^{m \times m}}^{(\alpha)}, \ j = \alpha + 1, \dots, k - 1 \\
    \br{y_{\alpha}^{m \times 1}}^{(\alpha + 1)}             & = L \br{y_{\alpha}^{m \times 1}}^{(\alpha)} \\
    \br{A_{\alpha, k}^{m \times l}}^{(\alpha + 1)}          & = L \br{A_{\alpha, k}^{m \times l}}^{(\alpha)}, \text{ если } l \neq 0
\end{align*}
Обновляем блоки, стоящие ниже строки выбранного блока и столбец свободных членов
\begin{align*}
    \br{A_{i, j}^{m \times m}}^{(\alpha + 1)} & = \br{A_{i, j}^{m \times m}}^{(\alpha)} - \br{A_{i, \alpha}^{m \times m}}^{(\alpha)} \br{A_{\alpha, j}^{m \times m}}^{(\alpha + 1)}, \ i, j = \alpha + 1, \dots, k - 1 \\
    \br{A_{i, k}^{m \times l}}^{(\alpha + 1)} & = \br{A_{i, k}^{m \times l}}^{(\alpha)} - \br{A_{i, \alpha}^{m \times m}}^{(\alpha)} \br{A_{\alpha, k}^{m \times l}}^{(\alpha + 1)}, \ i = \alpha + 1, \dots k - 1, \ l \neq 0 \\
    \br{y_{i}^{m \times 1}}^{(\alpha + 1)} & = \br{y_{i}^{m \times 1}}^{(\alpha)} - \br{A_{i, \alpha}^{m \times m}}^{(\alpha)} \br{y_{\alpha}^{m \times 1}}^{(\alpha + 1)}, \ i = \alpha + 1, \dots, k - 1 \\
    \br{A_{k, j}^{l \times m}}^{(\alpha + 1)} & = \br{A_{k, \alpha}^{l \times m}}^{(\alpha)} - \br{A_{k, j}^{l \times m}}^{(\alpha)} \br{A_{\alpha, j}^{m \times m}}^{(\alpha + 1)}, \ j = \alpha + 1, \dots, k - 1, \ l \neq 0 \\
    \br{A_{k, k}^{l \times l}}^{(\alpha + 1)} & = \br{A_{k, k}^{l \times l}}^{(\alpha)} - \br{A_{k, \alpha}^{l \times m}}^{(\alpha)} \br{A_{\alpha, k}^{m \times l}}^{(\alpha + 1)}, \ l \neq 0 \\
    \br{y_{k}^{l \times 1}}^{(\alpha + 1)} & = \br{y_{k}^{l \times 1}}^{(\alpha)} - \br{A_{k, \alpha}^{l \times m}}^{(\alpha)} \br{y_{\alpha}^{m \times 1}}^{(\alpha + 1)}, \ l \neq 0 \\
    \br{A_{i, \alpha}^{m \times m}}^{(\alpha + 1)} & = 0, \ i = \alpha + 1, \dots, k - 1 \\
    \br{A_{k, \alpha}^{l \times m}}^{(\alpha + 1)} & = 0, \ l \neq 0 \\
\end{align*}

После выполнения всех шагов получим матрицу
\[
    \begin{pmatrix}
        E^{m \times m} & \br{A_{0, 1}^{m \times m}}^{(k)}     & \br{A_{0, 2}^{m \times m}}^{(k)}     & \dots  & \br{A_{0, k - 1}^{m \times m}}^{(k)}     & \br{A_{0, k}^{m \times l}}^{(k)}        & \br{y_{0}^{m \times 1}}^{(k)}     \\
        0              & E^{m \times m}                       & \br{A_{1, 2}^{m \times m}}^{(k)}     & \dots  & \br{A_{1, k - 1}^{m \times m}}^{(k)}     & \br{A_{1, k}^{m \times l}}^{(k)}        & \br{y_{1}^{m \times 1}}^{(k)}     \\
        \vdots         & \vdots                               & \vdots                               & \ddots & \vdots                                   & \vdots                                  & \vdots                            \\
        0              & 0                                    &  0                                   & \dots  & E^{m \times m}                           & \br{A_{k - 1, k}^{m \times l}}^{\br{k}} & \br{y_{k - 1}^{m \times 1}}^{(k)} \\
        0              & 0                                    &  0                                   & \dots  & 0                                        & \br{A_{k, k}^{l \times l}}^{(k)}        & \br{y_{k}^{l \times 1}}^{(k)}
    \end{pmatrix}.
\]

Если $l \neq 0$, то остается решить систему из последней строчки. Для этого будет необходимо найти обратную матрицу размера $l \times l$.
\begin{align*}
    \br{{A_{k, k}^{l \times l}}}^{(k + 1)} & = E^{l \times l} \\
    \br{{y_{k}^{l \times 1}}}^{(k + 1)} & = \br{\br{{A_{k, k}^{l \times l}}}^{(k)}}^{-1} \br{{y_{k}^{l \times 1}}}^{(k)}
\end{align*}

\subsection{Обратный ход метода Гаусса}
После прямого хода диагональные блоки единичные. Индексы хода опускаем. Находим неизвестные в текущем порядке столбцов, начиная с последней блочной строки:
\begin{align*}
    x_{k}^{l \times 1} & = y_{k}^{l \times 1}, \ l \neq 0, \\
    x_{i}^{m \times 1} & = y_{i}^{m \times 1} - \sum_{j = i + 1}^{k - 1} A_{i, j}^{m \times m} x_{j}^{m \times 1} - A_{i, k}^{m \times l} x_{k}^{l \times 1}, \ i = k - 1, \dots, 0, l \neq 0, \\
    x_{i}^{m \times 1} & = y_{i}^{m \times 1} - \sum_{j = i + 1}^{k - 1} A_{i, j}^{m \times m} x_{j}^{m \times 1}, \ i = k - 1, \dots, 0, \ l = 0.
\end{align*}

Для восстановления исходного порядка неизвестных применяем к блокам вектора $x$ сохраненные перестановки столбцов в обратном порядке.

% \section{Число операций}
%     Оценка дана для невырожденной матрицы. 
%     На шаге $\alpha$ прямого хода вычисление множителей требует $n - \alpha - 1$ делений, обновление правой части - $2(n - \alpha - 1)$ операций, обновление матрицы - $2(n - \alpha - 1)^2$ операций. Число операций прямого хода равно
% \[
%     \sum_{\alpha = 0}^{n - 1}\bigl(2(n - \alpha - 1)^2 + 3(n - \alpha - 1)\bigr)
%     = \frac{2}{3} n^3 + \frac{1}{2} n^2 - \frac{7}{6} n.
% \]
% Обратный ход требует
% \[
%     \sum_{i = 0}^{n - 1}\bigl(2(n - i - 1) + 1\bigr) = n^2.
% \]
% Суммарное число арифметических операций:
% \[
%     \frac{2}{3} n^3 + \frac{3}{2} n^2 - \frac{7}{6} n.
% \]

% Поиск главного элемента на шаге $\alpha$ требует $(n - \alpha)^2 - 1$ сравнений одной проверки максимума на ноль. Число сравнений равно
% \[
%     \sum_{\alpha = 0}^{n - 1}(n - \alpha)^2 = \frac{1}{3} n^3 + \frac{1}{2} n^2 + \frac{1}{6} n.
% \]
% Общее число учитываемых операций равно
% \[
%     S(n, m) = n^3 + 2 n^2 - n
%     = C_1 n^3 + C_2 n^2 m + C_3 n m^2 + C_4 m^3 + O(n^2 + n m + m^2),
% \]
% \[
%     \boxed{C_1 = 1, \qquad C_2 = C_3 = C_4 = 0.}
% \]

% \section{Проверка формул}
% При $m = 1$. Формулы принимают обычный вид:
% \[
%     a_{i, j}^{(\alpha + 1)} = a_{i, j}^{(\alpha)} -
%     \frac{a_{i, \alpha}^{(\alpha)}}{a_{\alpha \alpha}^{(\alpha)}} a_{\alpha, j}^{(\alpha)},
%     \ \ \ \ 
%     y_i^{(\alpha + 1)} = y_i^{(\alpha)} -
%     \frac{a_{i, \alpha}^{(\alpha)}}{a_{\alpha \alpha}^{(\alpha)}} y_\alpha^{(\alpha)},
%     \ \ \ \ i, j > \alpha.
% \]
% В обратном ходе сумма внутри диагонального блока пуста, поэтому
% \[
%     x_i^{(n)} = \frac{y_i^{(n)} - \sum_{j = i + 1}^{n - 1} a_{i, j}^{(n)} x_j^{(n)}}{a_{i, i}^{(n)}},
%     \ \ \ \ \  i = n - 1, \ldots, 0.
% \]

% При $m = n$ Поиск идёт по $p, q \geq \alpha$, исключение - по $p, q > \alpha$, что даёт те же скалярные формулы. В обратном ходе сумма по блокам пуста, а сумма внутри блока даёт обычную формулу обратной подстановки.

% В обоих случаях число операций с учётом сравнений совпадает с неблочной реализацией:
% \[
%     S(n, 1) = S(n, n) = n^3 + 2 n^2 - n.
% \]
% Сложность остаётся $O(n^3)$.

\end{document}
